% =====================================================================
%  Panel captions for "The Categorical Ladder"
%  Generated by generate_panels.py -- do not edit by hand.
%  Usage:  % =====================================================================
%  Panel captions for "The Categorical Ladder"
%  Generated by generate_panels.py -- do not edit by hand.
%  Usage:  % =====================================================================
%  Panel captions for "The Categorical Ladder"
%  Generated by generate_panels.py -- do not edit by hand.
%  Usage:  % =====================================================================
%  Panel captions for "The Categorical Ladder"
%  Generated by generate_panels.py -- do not edit by hand.
%  Usage:  \input{figures/ladder-captions}  then  \captionPanelOne  etc.
% =====================================================================

\newcommand{\captionPanelOne}{\caption{\textbf{The closed-ladder
invariants, and an honest account of what each one buys.}
(\textbf{A})~Circulation $\varrho=-\sum_i\log(1-\pi_i)$ of a uniform cycle
over cycle length $n$ and rung power $\pi$. The surface rises without
bound in $\pi$ and linearly in $n$: a cycle deposits residue on every
circuit, which is why a closed ladder is not a no-op even though it
returns to the state it started from.
(\textbf{B})~Both invariants evaluated on all eight rotations of one
random eight-rung cycle. Each is exactly flat; measured over $2000$ random
cycles the rotation deviation is $3.6\times10^{-15}$ for $\varrho$ and
$4.4\times10^{-16}$ for $u$, so both are rotation invariants and a closed
ladder is well defined up to rotation.
(\textbf{C})~Circulation against composite power for $1200$ random linear
ladders, with the curve $-\log(1-\Pi)$ drawn through them. Every point
lies on the curve. This is the redundancy stated in the text rather than
concealed: for a \emph{linear} ladder $\varrho$ is a monotone
reparametrisation of composite power and adds no information whatever. Its
value is that it survives the passage to a cycle, where composite power is
undefined for want of a target.
(\textbf{D})~Absolute difference in uniformity between $4000$ pairs of
ladders constructed to share a composite power \emph{exactly}. Were $u$ a
function of composite power the histogram would collapse onto the red line
at zero; instead $69\%$ of matched pairs differ in $u$. This is the
measurement establishing that $u$ carries information composite power does
not, and it is the reason two invariants are carried rather than one.}}

\newcommand{\captionPanelTwo}{\caption{\textbf{A registered formulation of
aromaticity, refuted by cyclohexane, and the corrected two-invariant
statement.}
(\textbf{A})~Eight rings in $(\varrho/n,\,u,\,n)$ space, aromatic in red
and saturated in blue, with the separating plane at the midpoint of the
circulation margin. The classes separate along the circulation axis and
not along the uniformity axis.
(\textbf{B})~Uniformity alone, sorted by ring. We predicted before running
that a ring is aromatic exactly when its power profile is rotation
invariant, i.e.\ when $u=1$. Cyclohexane, cyclopentane and cyclobutane all
attain $u=1.000$ exactly and none is aromatic, so the prediction is
refuted by the second-most familiar ring in chemistry. The panel is
included because it is the refutation and not despite it.
(\textbf{C})~Circulation per rung, same ordering. The shaded band is the
empty interval between the largest saturated value ($0.400$) and the
smallest aromatic one ($0.693$), a margin of $+0.293$; no ring lies inside
it. Circulation is what separates the classes.
(\textbf{D})~The two coordinates together. Benzene and cyclohexane sit at
the same height ($u=1.000$) and are separated only horizontally;
benzene and pyridine sit at the same side of the threshold and are
separated only vertically. Neither invariant classifies alone and the pair
does, which is the corrected statement.}}

\newcommand{\captionPanelThree}{\caption{\textbf{Substitution is a letter
rewrite, and its effect depends on position rather than on count.}
(\textbf{A})~Loss of uniformity over the number of ring letters rewritten
and the power assigned to the substituted rung. The surface is identically
zero along the whole $k=0$ edge---a substitution outside the cycle cannot
move an invariant that is a function of the cycle's profile alone---and
rises once any ring letter is touched.
(\textbf{B})~The positional prediction on five cases. Peripheral
substitution gives $\Delta u=0.000$ exactly; every ring substitution gives
at least $0.086$. This is the asymmetry an earlier element-free matcher
measured and reported as its least convincing result, four cross-element
pairings across six molecule pairs; here it follows from where the
rewritten letter sits relative to the cycle rather than from a corpus.
(\textbf{C})~The cyclic power profiles themselves, drawn as step functions
around the six positions of the ring. Benzene is flat; each substitution
raises one position, and it is the resulting dispersion---not the pattern
of raised positions---that $u$ reads.
(\textbf{D})~A second registered expectation, also refuted. We predicted
$\Delta u$ would fall at three symmetric substitutions, since a symmetric
trisubstitution restores the rotational symmetry of the substitution
pattern (grey, dashed). It does not: the measured values rise monotonically
$0.086\to0.105\to0.107$, because $u$ measures dispersion of the profile's
values and symmetry of the pattern is not sameness of the values. An
invariant reading the symmetry group of the substitution would behave as
predicted; $u$ is not that invariant, and we record the gap rather than
redefine the quantity after seeing the result.}}

\newcommand{\captionPanelFour}{\caption{\textbf{A rung is inert exactly
when its step-pair repeats, which makes the halting condition derived
rather than declared.}
(\textbf{A})~Nine independently generated histories over a seven-state
set, each traced as the size of its determined set $\Delta$ against step
index. Red markers are steps whose unordered state-pair had already been
traversed. Every red marker falls on a plateau and every rise is
unmarked, which is the biconditional displayed one history at a time.
(\textbf{B})~The four possible outcomes counted over $2061$ steps. The two
green bars are the consistent cases, $1716$ fresh non-repetitions and
$345$ inert repetitions; the two red bars are the violations and both are
exactly zero. That both green bars are populated is the control: a run in
which no step ever repeated would satisfy the biconditional vacuously.
(\textbf{C})~Fourteen further histories showing the accumulation directly.
$\Delta$ is a union over steps taken and a union admits no deletion, so
the curves are monotone; the flat segments are exactly the repetitions.
(\textbf{D})~Fraction of steps that are inert against history length, over
$400$ histories at each length. Inert steps accumulate as the finite
supply of fresh state-pairs is exhausted, so a ladder run long enough
halts not by reaching a declared threshold but by running out of
non-repetitions.}}

\newcommand{\captionPanelFive}{\caption{\textbf{Why a coarser comparison
outperforms a finer one: refinement appends rungs that split classes while
depositing no circulation.}
(\textbf{A})~Two surfaces over refinement radius and ladder size. The flat
green plane at zero is the circulation gained by refinement; the red
wireframe rising above it is the number of label classes split. The gap
between them is the whole of the effect.
(\textbf{B})~Mean circulation gained per radius, which is zero to machine
precision at every radius. Refinement appends step-pairs already
traversed, and such steps are inert.
(\textbf{C})~Mean classes split per radius over the same ladders, growing
monotonically. The appended rungs do discriminate---they simply
discriminate without determining anything new.
(\textbf{D})~The measurement this explains, reproduced from the earlier
element-free matcher: mean class overlap for bioisosteric (close) and
unrelated (far) pairs against refinement radius, with the separation
margin as bars. The margin is $+0.375$ at radius $0$ and exactly $0.000$
at radii $1$, $2$ and $3$. Panels (\textbf{B}) and (\textbf{C}) supply a
mechanism for panel (\textbf{D}); they do not re-measure it, and the text
marks this category as a consistency check rather than as independent
confirmation.}}

\newcommand{\captionPanelSix}{\caption{\textbf{A correction to the
sharpest claim of both prior ladder papers: the counter-intuitive
direction is a property of the parametrisation, not of the ladder.}
(\textbf{A})~The additive sensitivity $\partial\Pi/\partial\pi_j =
P/(1-\pi_j)$ over the rung's own power and the residual fraction $P$. The
surface is increasing in $\pi_j$, which is the basis of the prior claim
that control lies at the strongest rung.
(\textbf{B})~The same ladder priced two ways. Under an additive increment,
gain differs across rungs and is largest at the strongest; under an
increment proportional to a rung's own remaining headroom---which is what
``improve this by ten percent'' ordinarily means---the factors $(1-\pi_j)$
cancel and the gain is $\delta P$ at every rung.
(\textbf{C})~Over $3000$ random ladders the additive argmax is the
strongest rung in every case, reproducing the prior result exactly; under
the proportional parametrisation there is no argmax to report because the
sensitivities are equal.
(\textbf{D})~Spread of sensitivity across rungs. The additive distribution
is broad; the proportional spread is a single line at zero, measured at
$1.1\times10^{-16}$ over $5000$ ladders. Both prior papers validated their
analytic derivative against a numerical one, a check that cannot
distinguish these two cases because both of its terms are the same
additive quantity.}}

\newcommand{\captionPanelSeven}{\caption{\textbf{Carrier elimination and
the single hypothesis whose failure destroys it.}
(\textbf{A})~Readout over the transit lengths of two carriers realising
the same contact sequence. Under local effect the readout is the flat
green plane: geometry is invisible to it. With a carrier-dependent
perturbation the red wireframe tilts, and the two carriers no longer
agree.
(\textbf{B})~Three counts over $3000$ trials. Carriers realising the same
sequence agree in every trial; but that row is true by construction once a
readout is defined to consult the terminal state alone, and is reported as
a consistency check. The load-bearing rows are the two controls: genuinely
different sequences separate, and a non-local contact effect separates.
Without them the first row would be consistent with a readout that ignores
its input entirely.
(\textbf{C})~Readout against carrier transit length at four violation
strengths. The horizontal green line is the eliminable case. Each coloured
curve is a locality violation, and the violation appears as a
\emph{slope}---a dependence of the reported value on a property of the
carrier that the theorem says cannot be recovered.
(\textbf{D})~Carrier-dependent deviation against violation strength
$\kappa$. The deviation grows from exactly zero and vanishes as
$\kappa\to0$, which is the diagnostic: in mass spectrometry this is a
dependence of measured mass on total ion current, and in catalysis a
dependence of a step's effect on which other sites are occupied. These are
the same violation under two names.}}
  then  \captionPanelOne  etc.
% =====================================================================

\newcommand{\captionPanelOne}{\caption{\textbf{The closed-ladder
invariants, and an honest account of what each one buys.}
(\textbf{A})~Circulation $\varrho=-\sum_i\log(1-\pi_i)$ of a uniform cycle
over cycle length $n$ and rung power $\pi$. The surface rises without
bound in $\pi$ and linearly in $n$: a cycle deposits residue on every
circuit, which is why a closed ladder is not a no-op even though it
returns to the state it started from.
(\textbf{B})~Both invariants evaluated on all eight rotations of one
random eight-rung cycle. Each is exactly flat; measured over $2000$ random
cycles the rotation deviation is $3.6\times10^{-15}$ for $\varrho$ and
$4.4\times10^{-16}$ for $u$, so both are rotation invariants and a closed
ladder is well defined up to rotation.
(\textbf{C})~Circulation against composite power for $1200$ random linear
ladders, with the curve $-\log(1-\Pi)$ drawn through them. Every point
lies on the curve. This is the redundancy stated in the text rather than
concealed: for a \emph{linear} ladder $\varrho$ is a monotone
reparametrisation of composite power and adds no information whatever. Its
value is that it survives the passage to a cycle, where composite power is
undefined for want of a target.
(\textbf{D})~Absolute difference in uniformity between $4000$ pairs of
ladders constructed to share a composite power \emph{exactly}. Were $u$ a
function of composite power the histogram would collapse onto the red line
at zero; instead $69\%$ of matched pairs differ in $u$. This is the
measurement establishing that $u$ carries information composite power does
not, and it is the reason two invariants are carried rather than one.}}

\newcommand{\captionPanelTwo}{\caption{\textbf{A registered formulation of
aromaticity, refuted by cyclohexane, and the corrected two-invariant
statement.}
(\textbf{A})~Eight rings in $(\varrho/n,\,u,\,n)$ space, aromatic in red
and saturated in blue, with the separating plane at the midpoint of the
circulation margin. The classes separate along the circulation axis and
not along the uniformity axis.
(\textbf{B})~Uniformity alone, sorted by ring. We predicted before running
that a ring is aromatic exactly when its power profile is rotation
invariant, i.e.\ when $u=1$. Cyclohexane, cyclopentane and cyclobutane all
attain $u=1.000$ exactly and none is aromatic, so the prediction is
refuted by the second-most familiar ring in chemistry. The panel is
included because it is the refutation and not despite it.
(\textbf{C})~Circulation per rung, same ordering. The shaded band is the
empty interval between the largest saturated value ($0.400$) and the
smallest aromatic one ($0.693$), a margin of $+0.293$; no ring lies inside
it. Circulation is what separates the classes.
(\textbf{D})~The two coordinates together. Benzene and cyclohexane sit at
the same height ($u=1.000$) and are separated only horizontally;
benzene and pyridine sit at the same side of the threshold and are
separated only vertically. Neither invariant classifies alone and the pair
does, which is the corrected statement.}}

\newcommand{\captionPanelThree}{\caption{\textbf{Substitution is a letter
rewrite, and its effect depends on position rather than on count.}
(\textbf{A})~Loss of uniformity over the number of ring letters rewritten
and the power assigned to the substituted rung. The surface is identically
zero along the whole $k=0$ edge---a substitution outside the cycle cannot
move an invariant that is a function of the cycle's profile alone---and
rises once any ring letter is touched.
(\textbf{B})~The positional prediction on five cases. Peripheral
substitution gives $\Delta u=0.000$ exactly; every ring substitution gives
at least $0.086$. This is the asymmetry an earlier element-free matcher
measured and reported as its least convincing result, four cross-element
pairings across six molecule pairs; here it follows from where the
rewritten letter sits relative to the cycle rather than from a corpus.
(\textbf{C})~The cyclic power profiles themselves, drawn as step functions
around the six positions of the ring. Benzene is flat; each substitution
raises one position, and it is the resulting dispersion---not the pattern
of raised positions---that $u$ reads.
(\textbf{D})~A second registered expectation, also refuted. We predicted
$\Delta u$ would fall at three symmetric substitutions, since a symmetric
trisubstitution restores the rotational symmetry of the substitution
pattern (grey, dashed). It does not: the measured values rise monotonically
$0.086\to0.105\to0.107$, because $u$ measures dispersion of the profile's
values and symmetry of the pattern is not sameness of the values. An
invariant reading the symmetry group of the substitution would behave as
predicted; $u$ is not that invariant, and we record the gap rather than
redefine the quantity after seeing the result.}}

\newcommand{\captionPanelFour}{\caption{\textbf{A rung is inert exactly
when its step-pair repeats, which makes the halting condition derived
rather than declared.}
(\textbf{A})~Nine independently generated histories over a seven-state
set, each traced as the size of its determined set $\Delta$ against step
index. Red markers are steps whose unordered state-pair had already been
traversed. Every red marker falls on a plateau and every rise is
unmarked, which is the biconditional displayed one history at a time.
(\textbf{B})~The four possible outcomes counted over $2061$ steps. The two
green bars are the consistent cases, $1716$ fresh non-repetitions and
$345$ inert repetitions; the two red bars are the violations and both are
exactly zero. That both green bars are populated is the control: a run in
which no step ever repeated would satisfy the biconditional vacuously.
(\textbf{C})~Fourteen further histories showing the accumulation directly.
$\Delta$ is a union over steps taken and a union admits no deletion, so
the curves are monotone; the flat segments are exactly the repetitions.
(\textbf{D})~Fraction of steps that are inert against history length, over
$400$ histories at each length. Inert steps accumulate as the finite
supply of fresh state-pairs is exhausted, so a ladder run long enough
halts not by reaching a declared threshold but by running out of
non-repetitions.}}

\newcommand{\captionPanelFive}{\caption{\textbf{Why a coarser comparison
outperforms a finer one: refinement appends rungs that split classes while
depositing no circulation.}
(\textbf{A})~Two surfaces over refinement radius and ladder size. The flat
green plane at zero is the circulation gained by refinement; the red
wireframe rising above it is the number of label classes split. The gap
between them is the whole of the effect.
(\textbf{B})~Mean circulation gained per radius, which is zero to machine
precision at every radius. Refinement appends step-pairs already
traversed, and such steps are inert.
(\textbf{C})~Mean classes split per radius over the same ladders, growing
monotonically. The appended rungs do discriminate---they simply
discriminate without determining anything new.
(\textbf{D})~The measurement this explains, reproduced from the earlier
element-free matcher: mean class overlap for bioisosteric (close) and
unrelated (far) pairs against refinement radius, with the separation
margin as bars. The margin is $+0.375$ at radius $0$ and exactly $0.000$
at radii $1$, $2$ and $3$. Panels (\textbf{B}) and (\textbf{C}) supply a
mechanism for panel (\textbf{D}); they do not re-measure it, and the text
marks this category as a consistency check rather than as independent
confirmation.}}

\newcommand{\captionPanelSix}{\caption{\textbf{A correction to the
sharpest claim of both prior ladder papers: the counter-intuitive
direction is a property of the parametrisation, not of the ladder.}
(\textbf{A})~The additive sensitivity $\partial\Pi/\partial\pi_j =
P/(1-\pi_j)$ over the rung's own power and the residual fraction $P$. The
surface is increasing in $\pi_j$, which is the basis of the prior claim
that control lies at the strongest rung.
(\textbf{B})~The same ladder priced two ways. Under an additive increment,
gain differs across rungs and is largest at the strongest; under an
increment proportional to a rung's own remaining headroom---which is what
``improve this by ten percent'' ordinarily means---the factors $(1-\pi_j)$
cancel and the gain is $\delta P$ at every rung.
(\textbf{C})~Over $3000$ random ladders the additive argmax is the
strongest rung in every case, reproducing the prior result exactly; under
the proportional parametrisation there is no argmax to report because the
sensitivities are equal.
(\textbf{D})~Spread of sensitivity across rungs. The additive distribution
is broad; the proportional spread is a single line at zero, measured at
$1.1\times10^{-16}$ over $5000$ ladders. Both prior papers validated their
analytic derivative against a numerical one, a check that cannot
distinguish these two cases because both of its terms are the same
additive quantity.}}

\newcommand{\captionPanelSeven}{\caption{\textbf{Carrier elimination and
the single hypothesis whose failure destroys it.}
(\textbf{A})~Readout over the transit lengths of two carriers realising
the same contact sequence. Under local effect the readout is the flat
green plane: geometry is invisible to it. With a carrier-dependent
perturbation the red wireframe tilts, and the two carriers no longer
agree.
(\textbf{B})~Three counts over $3000$ trials. Carriers realising the same
sequence agree in every trial; but that row is true by construction once a
readout is defined to consult the terminal state alone, and is reported as
a consistency check. The load-bearing rows are the two controls: genuinely
different sequences separate, and a non-local contact effect separates.
Without them the first row would be consistent with a readout that ignores
its input entirely.
(\textbf{C})~Readout against carrier transit length at four violation
strengths. The horizontal green line is the eliminable case. Each coloured
curve is a locality violation, and the violation appears as a
\emph{slope}---a dependence of the reported value on a property of the
carrier that the theorem says cannot be recovered.
(\textbf{D})~Carrier-dependent deviation against violation strength
$\kappa$. The deviation grows from exactly zero and vanishes as
$\kappa\to0$, which is the diagnostic: in mass spectrometry this is a
dependence of measured mass on total ion current, and in catalysis a
dependence of a step's effect on which other sites are occupied. These are
the same violation under two names.}}
  then  \captionPanelOne  etc.
% =====================================================================

\newcommand{\captionPanelOne}{\caption{\textbf{The closed-ladder
invariants, and an honest account of what each one buys.}
(\textbf{A})~Circulation $\varrho=-\sum_i\log(1-\pi_i)$ of a uniform cycle
over cycle length $n$ and rung power $\pi$. The surface rises without
bound in $\pi$ and linearly in $n$: a cycle deposits residue on every
circuit, which is why a closed ladder is not a no-op even though it
returns to the state it started from.
(\textbf{B})~Both invariants evaluated on all eight rotations of one
random eight-rung cycle. Each is exactly flat; measured over $2000$ random
cycles the rotation deviation is $3.6\times10^{-15}$ for $\varrho$ and
$4.4\times10^{-16}$ for $u$, so both are rotation invariants and a closed
ladder is well defined up to rotation.
(\textbf{C})~Circulation against composite power for $1200$ random linear
ladders, with the curve $-\log(1-\Pi)$ drawn through them. Every point
lies on the curve. This is the redundancy stated in the text rather than
concealed: for a \emph{linear} ladder $\varrho$ is a monotone
reparametrisation of composite power and adds no information whatever. Its
value is that it survives the passage to a cycle, where composite power is
undefined for want of a target.
(\textbf{D})~Absolute difference in uniformity between $4000$ pairs of
ladders constructed to share a composite power \emph{exactly}. Were $u$ a
function of composite power the histogram would collapse onto the red line
at zero; instead $69\%$ of matched pairs differ in $u$. This is the
measurement establishing that $u$ carries information composite power does
not, and it is the reason two invariants are carried rather than one.}}

\newcommand{\captionPanelTwo}{\caption{\textbf{A registered formulation of
aromaticity, refuted by cyclohexane, and the corrected two-invariant
statement.}
(\textbf{A})~Eight rings in $(\varrho/n,\,u,\,n)$ space, aromatic in red
and saturated in blue, with the separating plane at the midpoint of the
circulation margin. The classes separate along the circulation axis and
not along the uniformity axis.
(\textbf{B})~Uniformity alone, sorted by ring. We predicted before running
that a ring is aromatic exactly when its power profile is rotation
invariant, i.e.\ when $u=1$. Cyclohexane, cyclopentane and cyclobutane all
attain $u=1.000$ exactly and none is aromatic, so the prediction is
refuted by the second-most familiar ring in chemistry. The panel is
included because it is the refutation and not despite it.
(\textbf{C})~Circulation per rung, same ordering. The shaded band is the
empty interval between the largest saturated value ($0.400$) and the
smallest aromatic one ($0.693$), a margin of $+0.293$; no ring lies inside
it. Circulation is what separates the classes.
(\textbf{D})~The two coordinates together. Benzene and cyclohexane sit at
the same height ($u=1.000$) and are separated only horizontally;
benzene and pyridine sit at the same side of the threshold and are
separated only vertically. Neither invariant classifies alone and the pair
does, which is the corrected statement.}}

\newcommand{\captionPanelThree}{\caption{\textbf{Substitution is a letter
rewrite, and its effect depends on position rather than on count.}
(\textbf{A})~Loss of uniformity over the number of ring letters rewritten
and the power assigned to the substituted rung. The surface is identically
zero along the whole $k=0$ edge---a substitution outside the cycle cannot
move an invariant that is a function of the cycle's profile alone---and
rises once any ring letter is touched.
(\textbf{B})~The positional prediction on five cases. Peripheral
substitution gives $\Delta u=0.000$ exactly; every ring substitution gives
at least $0.086$. This is the asymmetry an earlier element-free matcher
measured and reported as its least convincing result, four cross-element
pairings across six molecule pairs; here it follows from where the
rewritten letter sits relative to the cycle rather than from a corpus.
(\textbf{C})~The cyclic power profiles themselves, drawn as step functions
around the six positions of the ring. Benzene is flat; each substitution
raises one position, and it is the resulting dispersion---not the pattern
of raised positions---that $u$ reads.
(\textbf{D})~A second registered expectation, also refuted. We predicted
$\Delta u$ would fall at three symmetric substitutions, since a symmetric
trisubstitution restores the rotational symmetry of the substitution
pattern (grey, dashed). It does not: the measured values rise monotonically
$0.086\to0.105\to0.107$, because $u$ measures dispersion of the profile's
values and symmetry of the pattern is not sameness of the values. An
invariant reading the symmetry group of the substitution would behave as
predicted; $u$ is not that invariant, and we record the gap rather than
redefine the quantity after seeing the result.}}

\newcommand{\captionPanelFour}{\caption{\textbf{A rung is inert exactly
when its step-pair repeats, which makes the halting condition derived
rather than declared.}
(\textbf{A})~Nine independently generated histories over a seven-state
set, each traced as the size of its determined set $\Delta$ against step
index. Red markers are steps whose unordered state-pair had already been
traversed. Every red marker falls on a plateau and every rise is
unmarked, which is the biconditional displayed one history at a time.
(\textbf{B})~The four possible outcomes counted over $2061$ steps. The two
green bars are the consistent cases, $1716$ fresh non-repetitions and
$345$ inert repetitions; the two red bars are the violations and both are
exactly zero. That both green bars are populated is the control: a run in
which no step ever repeated would satisfy the biconditional vacuously.
(\textbf{C})~Fourteen further histories showing the accumulation directly.
$\Delta$ is a union over steps taken and a union admits no deletion, so
the curves are monotone; the flat segments are exactly the repetitions.
(\textbf{D})~Fraction of steps that are inert against history length, over
$400$ histories at each length. Inert steps accumulate as the finite
supply of fresh state-pairs is exhausted, so a ladder run long enough
halts not by reaching a declared threshold but by running out of
non-repetitions.}}

\newcommand{\captionPanelFive}{\caption{\textbf{Why a coarser comparison
outperforms a finer one: refinement appends rungs that split classes while
depositing no circulation.}
(\textbf{A})~Two surfaces over refinement radius and ladder size. The flat
green plane at zero is the circulation gained by refinement; the red
wireframe rising above it is the number of label classes split. The gap
between them is the whole of the effect.
(\textbf{B})~Mean circulation gained per radius, which is zero to machine
precision at every radius. Refinement appends step-pairs already
traversed, and such steps are inert.
(\textbf{C})~Mean classes split per radius over the same ladders, growing
monotonically. The appended rungs do discriminate---they simply
discriminate without determining anything new.
(\textbf{D})~The measurement this explains, reproduced from the earlier
element-free matcher: mean class overlap for bioisosteric (close) and
unrelated (far) pairs against refinement radius, with the separation
margin as bars. The margin is $+0.375$ at radius $0$ and exactly $0.000$
at radii $1$, $2$ and $3$. Panels (\textbf{B}) and (\textbf{C}) supply a
mechanism for panel (\textbf{D}); they do not re-measure it, and the text
marks this category as a consistency check rather than as independent
confirmation.}}

\newcommand{\captionPanelSix}{\caption{\textbf{A correction to the
sharpest claim of both prior ladder papers: the counter-intuitive
direction is a property of the parametrisation, not of the ladder.}
(\textbf{A})~The additive sensitivity $\partial\Pi/\partial\pi_j =
P/(1-\pi_j)$ over the rung's own power and the residual fraction $P$. The
surface is increasing in $\pi_j$, which is the basis of the prior claim
that control lies at the strongest rung.
(\textbf{B})~The same ladder priced two ways. Under an additive increment,
gain differs across rungs and is largest at the strongest; under an
increment proportional to a rung's own remaining headroom---which is what
``improve this by ten percent'' ordinarily means---the factors $(1-\pi_j)$
cancel and the gain is $\delta P$ at every rung.
(\textbf{C})~Over $3000$ random ladders the additive argmax is the
strongest rung in every case, reproducing the prior result exactly; under
the proportional parametrisation there is no argmax to report because the
sensitivities are equal.
(\textbf{D})~Spread of sensitivity across rungs. The additive distribution
is broad; the proportional spread is a single line at zero, measured at
$1.1\times10^{-16}$ over $5000$ ladders. Both prior papers validated their
analytic derivative against a numerical one, a check that cannot
distinguish these two cases because both of its terms are the same
additive quantity.}}

\newcommand{\captionPanelSeven}{\caption{\textbf{Carrier elimination and
the single hypothesis whose failure destroys it.}
(\textbf{A})~Readout over the transit lengths of two carriers realising
the same contact sequence. Under local effect the readout is the flat
green plane: geometry is invisible to it. With a carrier-dependent
perturbation the red wireframe tilts, and the two carriers no longer
agree.
(\textbf{B})~Three counts over $3000$ trials. Carriers realising the same
sequence agree in every trial; but that row is true by construction once a
readout is defined to consult the terminal state alone, and is reported as
a consistency check. The load-bearing rows are the two controls: genuinely
different sequences separate, and a non-local contact effect separates.
Without them the first row would be consistent with a readout that ignores
its input entirely.
(\textbf{C})~Readout against carrier transit length at four violation
strengths. The horizontal green line is the eliminable case. Each coloured
curve is a locality violation, and the violation appears as a
\emph{slope}---a dependence of the reported value on a property of the
carrier that the theorem says cannot be recovered.
(\textbf{D})~Carrier-dependent deviation against violation strength
$\kappa$. The deviation grows from exactly zero and vanishes as
$\kappa\to0$, which is the diagnostic: in mass spectrometry this is a
dependence of measured mass on total ion current, and in catalysis a
dependence of a step's effect on which other sites are occupied. These are
the same violation under two names.}}
  then  \captionPanelOne  etc.
% =====================================================================

\newcommand{\captionPanelOne}{\caption{\textbf{The closed-ladder
invariants, and an honest account of what each one buys.}
(\textbf{A})~Circulation $\varrho=-\sum_i\log(1-\pi_i)$ of a uniform cycle
over cycle length $n$ and rung power $\pi$. The surface rises without
bound in $\pi$ and linearly in $n$: a cycle deposits residue on every
circuit, which is why a closed ladder is not a no-op even though it
returns to the state it started from.
(\textbf{B})~Both invariants evaluated on all eight rotations of one
random eight-rung cycle. Each is exactly flat; measured over $2000$ random
cycles the rotation deviation is $3.6\times10^{-15}$ for $\varrho$ and
$4.4\times10^{-16}$ for $u$, so both are rotation invariants and a closed
ladder is well defined up to rotation.
(\textbf{C})~Circulation against composite power for $1200$ random linear
ladders, with the curve $-\log(1-\Pi)$ drawn through them. Every point
lies on the curve. This is the redundancy stated in the text rather than
concealed: for a \emph{linear} ladder $\varrho$ is a monotone
reparametrisation of composite power and adds no information whatever. Its
value is that it survives the passage to a cycle, where composite power is
undefined for want of a target.
(\textbf{D})~Absolute difference in uniformity between $4000$ pairs of
ladders constructed to share a composite power \emph{exactly}. Were $u$ a
function of composite power the histogram would collapse onto the red line
at zero; instead $69\%$ of matched pairs differ in $u$. This is the
measurement establishing that $u$ carries information composite power does
not, and it is the reason two invariants are carried rather than one.}}

\newcommand{\captionPanelTwo}{\caption{\textbf{A registered formulation of
aromaticity, refuted by cyclohexane, and the corrected two-invariant
statement.}
(\textbf{A})~Eight rings in $(\varrho/n,\,u,\,n)$ space, aromatic in red
and saturated in blue, with the separating plane at the midpoint of the
circulation margin. The classes separate along the circulation axis and
not along the uniformity axis.
(\textbf{B})~Uniformity alone, sorted by ring. We predicted before running
that a ring is aromatic exactly when its power profile is rotation
invariant, i.e.\ when $u=1$. Cyclohexane, cyclopentane and cyclobutane all
attain $u=1.000$ exactly and none is aromatic, so the prediction is
refuted by the second-most familiar ring in chemistry. The panel is
included because it is the refutation and not despite it.
(\textbf{C})~Circulation per rung, same ordering. The shaded band is the
empty interval between the largest saturated value ($0.400$) and the
smallest aromatic one ($0.693$), a margin of $+0.293$; no ring lies inside
it. Circulation is what separates the classes.
(\textbf{D})~The two coordinates together. Benzene and cyclohexane sit at
the same height ($u=1.000$) and are separated only horizontally;
benzene and pyridine sit at the same side of the threshold and are
separated only vertically. Neither invariant classifies alone and the pair
does, which is the corrected statement.}}

\newcommand{\captionPanelThree}{\caption{\textbf{Substitution is a letter
rewrite, and its effect depends on position rather than on count.}
(\textbf{A})~Loss of uniformity over the number of ring letters rewritten
and the power assigned to the substituted rung. The surface is identically
zero along the whole $k=0$ edge---a substitution outside the cycle cannot
move an invariant that is a function of the cycle's profile alone---and
rises once any ring letter is touched.
(\textbf{B})~The positional prediction on five cases. Peripheral
substitution gives $\Delta u=0.000$ exactly; every ring substitution gives
at least $0.086$. This is the asymmetry an earlier element-free matcher
measured and reported as its least convincing result, four cross-element
pairings across six molecule pairs; here it follows from where the
rewritten letter sits relative to the cycle rather than from a corpus.
(\textbf{C})~The cyclic power profiles themselves, drawn as step functions
around the six positions of the ring. Benzene is flat; each substitution
raises one position, and it is the resulting dispersion---not the pattern
of raised positions---that $u$ reads.
(\textbf{D})~A second registered expectation, also refuted. We predicted
$\Delta u$ would fall at three symmetric substitutions, since a symmetric
trisubstitution restores the rotational symmetry of the substitution
pattern (grey, dashed). It does not: the measured values rise monotonically
$0.086\to0.105\to0.107$, because $u$ measures dispersion of the profile's
values and symmetry of the pattern is not sameness of the values. An
invariant reading the symmetry group of the substitution would behave as
predicted; $u$ is not that invariant, and we record the gap rather than
redefine the quantity after seeing the result.}}

\newcommand{\captionPanelFour}{\caption{\textbf{A rung is inert exactly
when its step-pair repeats, which makes the halting condition derived
rather than declared.}
(\textbf{A})~Nine independently generated histories over a seven-state
set, each traced as the size of its determined set $\Delta$ against step
index. Red markers are steps whose unordered state-pair had already been
traversed. Every red marker falls on a plateau and every rise is
unmarked, which is the biconditional displayed one history at a time.
(\textbf{B})~The four possible outcomes counted over $2061$ steps. The two
green bars are the consistent cases, $1716$ fresh non-repetitions and
$345$ inert repetitions; the two red bars are the violations and both are
exactly zero. That both green bars are populated is the control: a run in
which no step ever repeated would satisfy the biconditional vacuously.
(\textbf{C})~Fourteen further histories showing the accumulation directly.
$\Delta$ is a union over steps taken and a union admits no deletion, so
the curves are monotone; the flat segments are exactly the repetitions.
(\textbf{D})~Fraction of steps that are inert against history length, over
$400$ histories at each length. Inert steps accumulate as the finite
supply of fresh state-pairs is exhausted, so a ladder run long enough
halts not by reaching a declared threshold but by running out of
non-repetitions.}}

\newcommand{\captionPanelFive}{\caption{\textbf{Why a coarser comparison
outperforms a finer one: refinement appends rungs that split classes while
depositing no circulation.}
(\textbf{A})~Two surfaces over refinement radius and ladder size. The flat
green plane at zero is the circulation gained by refinement; the red
wireframe rising above it is the number of label classes split. The gap
between them is the whole of the effect.
(\textbf{B})~Mean circulation gained per radius, which is zero to machine
precision at every radius. Refinement appends step-pairs already
traversed, and such steps are inert.
(\textbf{C})~Mean classes split per radius over the same ladders, growing
monotonically. The appended rungs do discriminate---they simply
discriminate without determining anything new.
(\textbf{D})~The measurement this explains, reproduced from the earlier
element-free matcher: mean class overlap for bioisosteric (close) and
unrelated (far) pairs against refinement radius, with the separation
margin as bars. The margin is $+0.375$ at radius $0$ and exactly $0.000$
at radii $1$, $2$ and $3$. Panels (\textbf{B}) and (\textbf{C}) supply a
mechanism for panel (\textbf{D}); they do not re-measure it, and the text
marks this category as a consistency check rather than as independent
confirmation.}}

\newcommand{\captionPanelSix}{\caption{\textbf{A correction to the
sharpest claim of both prior ladder papers: the counter-intuitive
direction is a property of the parametrisation, not of the ladder.}
(\textbf{A})~The additive sensitivity $\partial\Pi/\partial\pi_j =
P/(1-\pi_j)$ over the rung's own power and the residual fraction $P$. The
surface is increasing in $\pi_j$, which is the basis of the prior claim
that control lies at the strongest rung.
(\textbf{B})~The same ladder priced two ways. Under an additive increment,
gain differs across rungs and is largest at the strongest; under an
increment proportional to a rung's own remaining headroom---which is what
``improve this by ten percent'' ordinarily means---the factors $(1-\pi_j)$
cancel and the gain is $\delta P$ at every rung.
(\textbf{C})~Over $3000$ random ladders the additive argmax is the
strongest rung in every case, reproducing the prior result exactly; under
the proportional parametrisation there is no argmax to report because the
sensitivities are equal.
(\textbf{D})~Spread of sensitivity across rungs. The additive distribution
is broad; the proportional spread is a single line at zero, measured at
$1.1\times10^{-16}$ over $5000$ ladders. Both prior papers validated their
analytic derivative against a numerical one, a check that cannot
distinguish these two cases because both of its terms are the same
additive quantity.}}

\newcommand{\captionPanelSeven}{\caption{\textbf{Carrier elimination and
the single hypothesis whose failure destroys it.}
(\textbf{A})~Readout over the transit lengths of two carriers realising
the same contact sequence. Under local effect the readout is the flat
green plane: geometry is invisible to it. With a carrier-dependent
perturbation the red wireframe tilts, and the two carriers no longer
agree.
(\textbf{B})~Three counts over $3000$ trials. Carriers realising the same
sequence agree in every trial; but that row is true by construction once a
readout is defined to consult the terminal state alone, and is reported as
a consistency check. The load-bearing rows are the two controls: genuinely
different sequences separate, and a non-local contact effect separates.
Without them the first row would be consistent with a readout that ignores
its input entirely.
(\textbf{C})~Readout against carrier transit length at four violation
strengths. The horizontal green line is the eliminable case. Each coloured
curve is a locality violation, and the violation appears as a
\emph{slope}---a dependence of the reported value on a property of the
carrier that the theorem says cannot be recovered.
(\textbf{D})~Carrier-dependent deviation against violation strength
$\kappa$. The deviation grows from exactly zero and vanishes as
$\kappa\to0$, which is the diagnostic: in mass spectrometry this is a
dependence of measured mass on total ion current, and in catalysis a
dependence of a step's effect on which other sites are occupied. These are
the same violation under two names.}}
